% Part: lambda-calculus
% Chapter: syntax
% Section: terms

\documentclass[../../../include/open-logic-section]{subfiles}

\begin{document}

\olfileid{lam}{syn}{trm}
\olsection{పదాలు}

లాంబ్డా కలనశాస్త్రపు పదాలను అనంతమైన చరాల జాబితా $\Obj{v_0}$,
$\Obj{v_1}$, \dots, ``$\lambd$'' అనే సంకేతం, కుండలీకరణ చిహ్నాల
నుంచి ఆగమన పద్ధతిలో నిర్మిస్తాం. చరాలను సూచించడానికి $x$, $y$, $z$,
\dots{}లను, పదాలను సూచించడానికి $M$, $N$, $P$, \dots{}లను వాడతాం.

\begin{defn}[పదాలు] \ollabel{defn:term}
లాంబ్డా కలనశాస్త్రపు \emph{పదాల} సమితిని ఆగమన పద్ధతిలో కింది విధంగా నిర్వచిస్తాం:
\begin{enumerate}
  \item \ollabel{defn:term-var} $x$ ఒక చరం అయితే, $x$ ఒక పదం.
  \item \ollabel{defn:term-abs} $x$ ఒక చరమూ, $M$ ఒక పదమూ అయితే,
    $(\lambd[x][M])$ ఒక పదం.
  \item \ollabel{defn:term-app} $M$, $N$ రెండూ పదాలైతే,
    $(MN)$ ఒక పదం.
\end{enumerate}
\end{defn}

$(\lambd[x][M])$ అనే పదాన్ని \olref{defn:term-abs} ప్రకారం నిర్మిస్తే
అది \emph{అమూర్తీకరణ} ఫలితం. $\lambd[x]$లోని $x$ను
\emph{\tetoken{పరామితి}{!!{parameter}}} అంటాం.
\olref{defn:term-app} ప్రకారం ఏర్పడిన పదం $(MN)$ ఒక
\emph{ప్రయోగం} ఫలితం.

పైన ఇచ్చిన నిర్వచనం ప్రకారం పదాలను అన్ని కుండలీకరణలతో రాస్తాం.
పదం $(\lambd[x][((\lambd[x][x])(\lambd[x][(xx)]))])$ చూపినట్లు, ఇలా
రాయడం కొంత క్లిష్టమవుతుంది.
కుండలీకరణలను వదిలివేయడానికి కొన్ని సంప్రదాయాలను ప్రవేశపెడతాం.
అయితే \olref{defn:term}లో ఇచ్చిన అధికారిక నిర్వచనం వల్ల ఒక పదాన్ని
ఎలా నిర్మించారో సులభంగా గుర్తించవచ్చు. ఉదాహరణకు,
$(\lambd[x][((\lambd[x][x])(\lambd[x][(xx)]))])$ నిర్మాణంలోని చివరి
దశ, \tetoken{పరామితి}{!!{parameter}}~$x$ గల అమూర్తీకరణ అయి ఉండాలి.
అది $((\lambd[x][x])(\lambd[x][(xx)]))$ అనే పదంపై అమూర్తీకరణ
చేయడం వల్ల వస్తుంది. ఆ పదం రెండు పదాలను ప్రయోగించిన ఫలితం. ఆ రెండింటిలో ప్రతి
పదమూ అమూర్తీకరణ ఫలితమే. ఇలాగే మిగిలిన దశలనూ గుర్తించవచ్చు.

\begin{prob}
$(\lambd[g][(\lambd[x][(g (x x))])
  (\lambd[x][(g (x x))])])$ ఎలా నిర్మించబడిందో వివరించండి.
\end{prob}

\end{document}
